% Chapter 7 file (Pass 2 unification, D1): Floquet + Weighted/Echoes + Coboundary,
% in verbatim root-input order. No content reordered, none cut.
% The Weighted/Echoes section below is Chapter 8 material embedded here so that
% the reading order (Floquet, Weighted, Coboundary) is preserved exactly.

\section{The arithmetic Floquet normal form and critical recurrence law}\label{sec:floquet}
Put \[
 \alpha=1+\sqrt2,\quad \lambda=\alpha^2,\quad
 T=\begin{pmatrix}3&2\\4&3\end{pmatrix}.
\] Let K tend to infinity through odd primes, and define positive integers N,F by \[
 N+F\sqrt2=\alpha^K,\qquad N^2-2F^2=-1.
\] The moduli N may be composite. On the odd subspace of the standard N-dimensional torus quantization, let \(U_{N}\) be the phase-normalized quantization of T described in Subsection~\ref{sec:floquet}.1, and let \[
 B_N f(x)=2\cos(2\pi x/N)f(x),\quad
 U_{\varepsilon,N}=e^{i\varepsilon B_N}U_{N,\mathrm{Fl}},\quad
 d_N=(N-1)/2.
\]

\begin{theorem}[critical resonance law]\label{pass1:ch07_floquet_resonance-1} For every fixed real \(\tau\) and every integer m, with \(\varepsilon=\tau/\sqrt{2K}\), \[
 \lim_{K\longrightarrow\infty}
 \frac{1}{d_N}\operatorname{Tr}_{-}
       U_{\varepsilon,N}^{mK}
       =\exp(-m^2\tau^2/2).                                      \tag{FN.A}
\] Consequently the empirical eigenvalue measure of \(U_{\varepsilon,N}^K\) converges weakly on the unit circle to the wrapped normal law. For \(\tau\) nonzero its angular density on {[}-\(\pi\),\(\pi\)) is \[
 w_\tau(\theta)=\frac{1}{\sqrt{2\pi}|\tau|}
       \sum_{\ell\in\mathbb Z}
       \exp\!\left[-\frac{(\theta+2\pi\ell)^2}{2\tau^2}\right].   \tag{FN.B}
\] For \(\tau=0\) the limit is the point mass at 1. In particular the normalized resonance echo obeys \[
 d_N^{-2}\left|\operatorname{Tr}_{-}U_{\varepsilon,N}^{mK}\right|^2
       \longrightarrow e^{-m^2\tau^2}.                          \tag{FN.C}
\] This is a deterministic limit: there is no parameter average in (FN.A).
\end{theorem}

The principal finite estimate, proved before any limit is taken, is \[
 \left\|V_\varepsilon U_{\varepsilon,N}V_\varepsilon^*
            -e^{i\varepsilon\overline B}U_{N,\mathrm{Fl}}\right\|_{2,-}
       \leq \frac{12\varepsilon^2}{\sqrt K},\qquad K\geq11,        \tag{FN.D}
\] for an explicit unitary \(V_{\varepsilon}\). Here the norm is normalized Hilbert--Schmidt norm on the entire odd subspace. It does not control an individual level or an operator-norm spectral gap. The finite comparison error in (FN.A) is at most \(6|m|\tau^2/\sqrt K\); an additional Gaussian characteristic-function convergence error is \(o(1)\), not assigned that rate.

\subsection{Arithmetic and the fixed classical observable}\label{floquet:arithmetic-and-the-fixed-classical-observable}

The recurrence \((N_j,F_j)=(N_{j-1}+2F_{j-1},N_{j-1}+F_{j-1})\), starting at \((1,0)\), shows that \(N\) is odd. In \(\mathbb Z/N\mathbb Z\), put \[
 s=F^{-1},\qquad s^2=2,\qquad h=3+2s,\qquad q=-F.
\] Since \((1+s)^K=N+Fs=1\), we have \(h^K=1\). Also \(h-1=2(1+s)\) is a unit. Modulo every prime \(\ell\mid N\), the order of \(h\) is precisely \(K\): it divides the prime \(K\) and is not one. Hence \(1-h^j\) is a unit for \(1\leq j<K\), and \(1+h^j\) is a unit for every \(j\), because \(K\) is odd. In particular, \(\ell\equiv1\pmod K\) and \(N\equiv1\pmod{2K}\). Neither squarefreeness nor primality of \(N\) is required.

Use the counting-measure inner product and centered Weyl operators \[
 W(a,b)f(x)=e_N(a(x-b/2))f(x-b),\qquad
 S_u f(x)=f(u^{-1}x),\qquad Rf(x)=f(-x).
\] Divisions in Weyl exponents mean inverses in \(\mathbb Z/N\mathbb Z\). The multiplication and covariance formulas are \[
 W(a,b)W(c,d)=e_N((ad-bc)/2)W(a+c,b+d),
\] \[
 S_uW(a,b)S_u^*=W(u^{-1}a,ub),\qquad RW(v)R=W(-v).
\] Individual Weyl translations need not preserve parity; their symmetrized sums do. On the odd subspace take \[
 U_0=S_h^{-1},\qquad B=W(q,1)+W(q,1)^*,\qquad
 U_\varepsilon=e^{i\varepsilon B}U_0.                            \tag{FN.1}
\] In these coordinates \(U_0^K=I\) exactly.

For completeness, this is unitarily equivalent to the fixed-map formulation of Theorem~\ref{pass1:ch07_floquet_resonance-1}, including the kick, not merely its compressed observable. Let \(C_tf(x)=e_N(tx^2/2)f(x)\), and let \(F_N\) have kernel \(e_N(xy)/\sqrt N\). Set \[
 M=\begin{pmatrix}-1/(2s)&1\\-1/2&-s\end{pmatrix},\qquad
 V=C_{s^{-1}}S_{1/2}F_N C_{2s}.
\] Its factors have Weyl covariance matrices respectively \[
 \begin{pmatrix}1&t\\0&1\end{pmatrix},\quad
 \begin{pmatrix}u^{-1}&0\\0&u\end{pmatrix},\quad
 \begin{pmatrix}0&1\\-1&0\end{pmatrix}.
\] Multiplication gives covariance \(M\), \(\det M=1\), and \(T=M\operatorname{diag}(h,h^{-1})M^{-1}\). Define \(U_{N,\mathrm{Fl}}=VU_0V^*\). The index of \(V^*W(1,0)V\) is \(M^{-1}(1,0)=(-s,1/2)\). Further conjugation by \(S_2\), which commutes with \(U_0\), changes this to \((-s/2,1)=(-F,1)\pmod N\). Thus \(S_2V^*\) conjugates the entire physical kicked propagator to (FN.1). These unitaries all preserve parity.

With the usual semiclassical convention \(e^{i\,\mathrm{action}/\hbar}\), the classical kick amplitude associated with \(\varepsilon\) is proportional to \(\hbar\varepsilon\). Here \(\hbar\asymp1/N\). At the critical scale the amplitude is of order \(1/(N\sqrt K)\), not a fixed classical perturbation.

\subsection{Two summable phase estimates}\label{floquet:two-summable-phase-estimates}

Write \(C(A)=U_0 A U_0^*\), and \[
 B_j=C^j(B)=W(v_j)+W(-v_j),\qquad v_j=(qh^j,h^{-j}),
 \qquad \overline B=K^{-1}\sum_{j=0}^{K-1}B_j.                    \tag{FN.2}
\] The conjugation C shifts the label cyclically, so \(C(\bar B)=\bar B\). Define \[
 a_j=\left|e_N(q(h^j-h^{-j}))-1\right|,\qquad
 c_j=\frac{K-1-2j}{2K},\qquad G=\sum_{j=0}^{K-1}c_jB_j.          \tag{FN.3}
\] Then \(a_0=0,\ a_{K-j}=a_j\). Direct subtraction of the coefficients gives the exact cohomological equation \[
 G-C(G)=B-\overline B.                                         \tag{FN.4}
\] The coefficients have zero mean. Importantly, they are small at the half-period at which the large commutator phases accumulate.

Let \(P_{t}\) denote the second Pell sequence, \(\alpha^t=Q_t+P_t\sqrt2\). For the signed representative \(\delta\) of j mod K one has \[
 q(h^j-h^{-j})=-2P_{2\delta}\pmod N.                            \tag{FN.5}
\] Indeed \(h^\delta=(3+2s)^\delta=Q_{2\delta}+P_{2\delta}s\), and \(Fs=1\). For positive \(\delta=(K-1)/2-\ell\), \[
 \frac{P_{2\delta}}{N}
 =\frac{\lambda^{-\ell}}{\sqrt2\,\alpha}
       \frac{1-\lambda^{-2\delta}}{1-\lambda^{-K}}.              \tag{FN.6}
\] Put \(r=\lambda^{-1}\), \(A=4\pi/(\sqrt2\,\alpha)\). Using \(|e^{it}-1|\leq|t|\) and summing both signed representatives proves \[
 L_K:=\sum_{j=0}^{K-1}a_j
 \leq\frac{2A}{(1-r)(1-\lambda^{-K})}
 =\frac{2\sqrt2\pi}{1-\lambda^{-K}}<9\quad(K\geq3).              \tag{FN.7}
\] There is a sharper weighted estimate. At j=(K-1)/2-ell the weights at j and K-j have absolute values ell/K and (ell+1)/K. Therefore \[
 \sum_j |c_j|a_j
 \leq\frac{A}{K(1-\lambda^{-K})}\sum_{\ell\geq0}(2\ell+1)r^\ell
 =\frac{2\pi}{K(1-\lambda^{-K})}.                              \tag{FN.8}
\] The final identity follows from \(A(1+r)/(1-r)^2=2\pi\). The estimates are analytic inequalities; the regression's decimal phase values are only finite checks of them.

\subsection{Quadratic lifting and orthogonal commutator supports}\label{floquet:quadratic-lifting-and-orthogonal-commutator-supports}

\begin{lemma}[lifting]\label{pass1:ch07_floquet_resonance-2} Fix m. For all sufficiently large K, every signed m-term reciprocal relation \[
 \sum_{\nu=1}^m\epsilon_\nu h^{j_\nu}
 =\sum_{\nu=1}^m\epsilon_\nu h^{-j_\nu}=0\pmod N,
 \qquad \epsilon_\nu\in\{-1,1\},                               \tag{FN.9}
\] can be cyclically rotated to a polynomial \(P\) with \(P(\lambda)=0\). Its coefficient \(\ell^1\)-norm is at most \(m\), and its degree is at most \(K-\lceil K/m\rceil\). For \(m=4\), this holds for every odd prime \(K\geq11\), and the relation is then coefficientwise balanced.
\end{lemma}

\textbf{Proof.} Cancel opposite signs at identical labels. If anything remains, cut the cyclic support at a largest gap and rotate. There are at most \(m\) support points, giving the stated degree bound. Both congruences in (FN.9) survive. Divide \(P\) by \(Z^2-6Z+1\) over the integers. Its remainder \(aZ+b\) vanishes at \(h\) and \(h^{-1}\); their difference \(4s\) is a unit. Thus \(N\mid a\) and \(N\mid b\). A nonzero quadratic norm of \(P(\lambda)\) would be an integer divisible by \(N^2\), but \[
 |P(\lambda)P(\lambda^{-1})|
 \leq m^2\lambda^{K-\lceil K/m\rceil}<N^2                       \tag{FN.10}
\] for sufficiently large K at fixed m. This is impossible, so \(P(\lambda)=0\). For \(m=4\), a sufficient condition is \(\lambda^{K/4}(1-\lambda^{-K})^2>64\). It holds for \(K\geq11\): \(\lambda>5\), \(5^{11/4}>80\), and \(1-5^{-11}>9/10\) give a strict lower bound of 64.8. A nonzero polynomial of coefficient \(\ell^1\)-norm at most four cannot vanish at \(\lambda>3\), since its leading term strictly dominates the sum of all lower terms. Thus in this case the polynomial was zero. End of proof.

It follows that the \(2K(K-1)\) vectors \[
 \epsilon v_t+\delta v_s,\qquad 0\leq t<s<K,\quad
 \epsilon,\delta\in\{-1,1\},                                  \tag{FN.11}
\] are nonzero and distinct for \(K\geq11\). Any collision would give a balanced four-letter relation, forcing the original unordered signed pair. The matrices \(W(v)\) form an orthonormal basis for the full normalized Hilbert--Schmidt inner product, also for composite odd \(N\), because \(\operatorname{Tr}W(v)=N\mathbf1_{v=0}\).

Each of the four Weyl coefficients in \([B_t,B_s]\) has modulus \(a_{s-t}\). Different unordered time pairs have disjoint support by (FN.11). Since \(c_t-c_s=(s-t)/K\), exact orthogonality gives \[
 \|[G,\overline B]\|_{2,N}^2
 =\frac4{K^4}\sum_{j=1}^{K-1}(K-j)j^2a_j^2
 =\frac2{K^3}\sum_{j=1}^{K-1}j(K-j)a_j^2
 \leq\frac9K.                                                  \tag{FN.12}
\] Here \(\|A\|_{2,N}^2=\operatorname{Tr}(A^*A)/N\). The middle identity uses \(a_{K-j}=a_j\); the last uses \(j(K-j)\leq K^2/4\), \(a_j^2\leq2a_j\), and (FN.7).

The operator norm needs no orthogonality. Four Weyl coefficients, (FN.7), and (FN.8) give \[
 \|[G,B]\|_{\rm op}<\frac{26}{K},\qquad
 \|[B,\overline B]\|_{\rm op}<\frac{36}{K}.                     \tag{FN.13}
\] For a parity-preserving A, positivity of A*A implies \[
 \|A\|_{2,-}\leq\sqrt{\frac{2N}{N-1}}\,\|A\|_{2,N}.             \tag{FN.14}
\] No orthogonality of Weyl operators restricted to the odd subspace is assumed.

\subsection{A finite normal form without a small generator}\label{floquet:a-finite-normal-form-without-a-small-generator}

For finite Hermitian A,D, real \(\varepsilon\), and any unitarily invariant norm, \[
 \|e^{i\varepsilon A}e^{i\varepsilon D}
           -e^{i\varepsilon(A+D)}\|
 \leq\frac{\varepsilon^2}{2}\|[A,D]\|.                          \tag{FN.15}
\] This is the usual two-factor commutator estimate, not a new general product formula. Differentiate \(e^{i(\varepsilon-t)(A+D)}e^{itA}e^{itD}\). After factoring out unitaries, the integrand is \(e^{itA}De^{-itA}-D\). Integrate its derivative once more; the triangular integration domain has area \(\varepsilon^2/2\). Negative \(\varepsilon\) follows by replacing \(A,D\) with \(-A,-D\). In particular, \(|\varepsilon|\|A\|\) need not be small.

Set \(V_\varepsilon=e^{-i\varepsilon G}\). By (FN.4), \[
 V_\varepsilon U_\varepsilon V_\varepsilon^*
 =e^{-i\varepsilon G}e^{i\varepsilon B}
       e^{i\varepsilon C(G)}U_0.
\] Apply (FN.15) first to \(-G,B\), then to \(B-G,C(G)\). Their total sum is \(\bar B\), and \([B-G,C(G)]=[B-G,\bar B]\). Equations (FN.12)--(FN.14) give \[
 \begin{split}
 \|V_\varepsilon U_\varepsilon V_\varepsilon^*
       -e^{i\varepsilon\overline B}U_0\|_{2,-}
 &\leq\frac{\varepsilon^2}{2}
   \left(\frac{62}{K}+3\sqrt{\frac{2N}{(N-1)K}}\right)\\
 &\leq\frac{12\varepsilon^2}{\sqrt K},\qquad K\geq11.
 \end{split}\tag{FN.16}
\] For the last step it even suffices to use \(N\geq7\) and \(31/\sqrt{11}+(3/2)\sqrt{7/3}<12\). This proves (FN.D) for all real \(\varepsilon\). G may have operator norm of order K. Expanding its exponential in powers of \(\varepsilon\) and requiring \(\varepsilon K\) small would incorrectly discard the critical regime. Estimate (FN.15) avoids that requirement.

Since \(\bar B\) commutes with \(U_{0}\) and \(U_0^K=I\), telescoping unitary powers in the same norm yields, for any integer m, \[
 \left|\tau_-(U_\varepsilon^{mK})
       -\tau_-(e^{im\varepsilon K\overline B})\right|
 \leq12|m|\varepsilon^2\sqrt K,\qquad
 \tau_-(A)=d_N^{-1}\operatorname{Tr}_{-}A.                      \tag{FN.17}
\] For negative m use adjoints. Cauchy--Schwarz for the normalized trace justifies the passage from the norm to traces.

\subsection{A self-contained global law for the averaged observable}\label{floquet:a-self-contained-global-law-for-the-averaged-observable}

The remaining input is established here directly on the whole odd subspace; no unproved claim about the old arithmetic-moment corpus is imported. Define \[
 X_K=\sqrt{K/2}\,\overline B=(2K)^{-1/2}\sum_{j=0}^{K-1}B_j.
\]

\begin{lemma}[Gaussian average]\label{pass1:ch07_floquet_resonance-3} For each fixed integer \(r\geq1\), \[
 \tau_-(X_K^{2r})=(2r-1)!!+O_r(K^{-1})+O_r(K^r/N),               \tag{FN.18}
\] and for each fixed odd m, eventually \[
 |\tau_-(X_K^m)|\leq (2K)^{m/2}/(N-1).                         \tag{FN.19}
\] The empirical measures of \(X_{K}\) converge weakly and in quadratic Wasserstein distance to the standard normal law.
\end{lemma}

\textbf{Proof.} In the full Weyl representation \(\operatorname{Tr}(RW(v))=1\) for every v: the unique fixed point \(x=-b/2\) of RW(a,b) contributes phase 1. Thus the normalized odd trace of a word \(e_{N}\)(\(\phi\))W(v) is \[
 e_N(\phi)\frac{N\mathbf1_{v=0}-1}{N-1}.                       \tag{FN.20}
\] Expanding \(X_K^m\), the noncentral part is bounded by \((2K)^{m/2}/(N-1)\). The central part is the weighted count of the reciprocal relations (FN.9), normalized by \([N/(N-1)](2K)^{-m/2}\).

Here is the count including structurally forced nonbalanced relations. For a polynomial of coefficient \(\ell^1\)-norm at most m vanishing at \(\lambda\), a gap \(d>2\log_\lambda m\) separates two independently vanishing pieces. Write \(P=P_{\rm low}+Z^{a+d}P_{\rm high}\), with \(\deg P_{\rm low}\leq a\) and nonnegative exponents in \(P_{\rm high}\). If \(P_{\rm high}(\lambda)\ne0\), its algebraic-integer norm is a nonzero integer. Since \(|P_{\rm high}(\lambda^{-1})|\leq m\), we have \(|P_{\rm high}(\lambda)|\geq1/m\). Cancellation with \(|P_{\rm low}(\lambda)|\leq m\lambda^a\) would force \(\lambda^d\leq m^2\), a contradiction.

Repeatedly splitting gives components with span bounded in terms of m, for example by \(D_m=(m-1)(\lceil2\log_\lambda m\rceil+1)\). Every nonzero component has at least three signed occurrences: two distinct positive powers cannot cancel with signs of magnitude one. Restore previously cancelled opposite pairs as two-occurrence components. A nonbalanced 2r-letter word therefore has at most r-1 components. For each of finitely many partitions of the 2r positions, signs and bounded relative offsets, each component has at most K translations on the cyclic time set. There are only \(O_r(K^{r-1})\) such words. There is no additional factor K from choosing a cut: the count is over the component translations themselves, and the largest-gap cut was used only to prove bounded width. Their normalized contribution is \(O_r(1/K)\).

For odd m, a lifted P is an integer multiple of \(Z^2-6Z+1\), so \(P(1)\) is divisible by -4. But \(P(1)\) is the sum of an odd number of signs. Thus no central word exists once Lemma~\ref{pass1:ch07_floquet_resonance-2} applies, proving (FN.19).

A balanced \(2r\)-word pairs identical time labels with opposite signs. Assignments with repeated pair labels cause only \(O_r(K^{r-1})\) overcount. Each matching of the \(2r\) positions has \(K^r\) pair-label choices and \(2^r\) orientations. Its Weyl phase is a product, over crossing pairs, of factors \(e_N(q(h^a-h^{-a}))\) or their conjugates. Summing the signed symplectic products in word order shows that disjoint or nested pairs cancel, while a crossing pair retains one full symplectic product.

Each crossing difference of independent uniform pair labels is uniform modulo \(K\). Independence of the different crossing factors is neither true in general nor needed. The inequality \(|\prod_l u_l-1|\leq\sum_l|u_l-1|\), for \(|u_l|=1\), and (FN.7) give average phase error at most nine times the number of crossings divided by \(K\). Hence each matching contributes \(1+O_r(1/K)\) after normalization. There are \((2r-1)!!\) matchings. Together with the nonbalanced count and (FN.20), this proves (FN.18).

The second-moment bound gives tightness. Higher even moments give uniform integrability when passing moments to subsequential limits. The Gaussian moment sequence is determinate, for example by Carleman\textquotesingle s criterion, since \(((2r-1)!!)^{1/(2r)}\leq\sqrt{2r}\). Every subsequential limit is standard normal. Convergence of the second moments upgrades weak convergence to quadratic Wasserstein convergence. End of proof.

\subsection{Critical and subcritical propagator spectra}\label{floquet:critical-and-subcritical-propagator-spectra}

For \(\varepsilon=\tau/\sqrt{2K}\), (FN.17) becomes \[
 \left|\tau_-(U_\varepsilon^{mK})-
              \tau_-(e^{im\tau X_K})\right|
       \leq \frac{6|m|\tau^2}{\sqrt K}.                        \tag{FN.21}
\] Lemma~\ref{pass1:ch07_floquet_resonance-3} and bounded continuity of exp(im \(\tau\) x) give (FN.A). Integer Fourier coefficients determine measures on the circle: trigonometric polynomials are dense in continuous functions there. Applying the characteristic function of a normal variable proves (FN.B), and taking absolute squares proves (FN.C). No eigenphase labelling across the old degenerate sectors is used or asserted.

There is also a genuine extension of the earlier subcritical global transfer, but for \textbf{folded phases} rather than individually isolated clusters. Let \(\varepsilon_{K}\) be nonzero and \(\tau_K=\varepsilon_K\sqrt{2K}\) tend to zero. If \(e^{i\theta_l}\) are the eigenvalues of \(U_{\varepsilon,K}\), define \[
 y_l=\frac{\operatorname{Arg}(e^{iK\theta_l})}{\tau_K},
 \qquad \operatorname{Arg}\in[-\pi,\pi).
\]

\begin{corollary}[subcritical folded Gaussian law]\label{pass1:ch07_floquet_resonance-4} The empirical measures of \(y_{l}\) converge in quadratic Wasserstein distance to the standard normal law whenever \(\varepsilon_K\sqrt K\) tends to zero.
\end{corollary}

\textbf{Proof with a quantitative tail control.} For two unitary matrices, the squared absolute overlaps of their orthonormal eigenbases form a doubly stochastic matrix. The corresponding average squared chord distance equals their normalized squared Hilbert--Schmidt difference. Decomposing this matrix into permutation matrices gives an eigenvalue matching with no larger average cost. The decomposition can be obtained by Hall's matching criterion and successive subtraction; it uses no eigenvalue simplicity. Thus (FN.16), telescoped K times, matches the powered spectrum to \(e^{i\tau_Kx_l}\), where \(x_{l}\) are the eigenvalues of \(X_{K}\), with root mean square chord error at most \(6\tau_K^2/\sqrt K\).

If \(|\tau_Kx_l|\leq\pi/2\), principal argument differences obey \(|\operatorname{Arg}u-\tau_Kx_l|\leq4|u-e^{i\tau_Kx_l}|\). On the complementary set, \(|\operatorname{Arg}u/\tau_K-x_l|\leq3|x_l|\). Consequently \[
 \begin{split}
 W_2(\operatorname{law}(y),\operatorname{law}(X_K))
 &\leq\frac{24|\tau_K|}{\sqrt K}
   +3\sqrt{\tau_-\!\left(X_K^2
       \mathbf1_{|X_K|>\pi/(2|\tau_K|)}\right)}\\
 &\leq\frac{24|\tau_K|}{\sqrt K}
       +\frac{6|\tau_K|}{\pi}\sqrt{\tau_-(X_K^4)}.
 \end{split}\tag{FN.22}
\] The fourth moment stays bounded by Lemma~\ref{pass1:ch07_floquet_resonance-3}, so the right side tends to zero. Combine with the \(W_{2}\) conclusion of that lemma. End of proof.

The former isolated-cluster argument required \(\varepsilon_KK^{3/2}\) tending to zero. Corollary~\ref{pass1:ch07_floquet_resonance-4} permits an additional factor K in the kick scale, but its observable is a folded spectral measure. It must not be restated as an individual-cluster rank or no-overlap theorem.

\subsection{A compact-test microscopic obstruction}\label{floquet:a-compact-test-microscopic-obstruction}

A negative statement about microscopic correlations needs less information than a positive universality theorem. The recurrence law above supplies such a statement through Fourier positivity. This section concerns the whole odd spectrum, with its multiplicities and the \textbf{global mean-density unfolding}. It is not a claim about a further symmetry sector or an unfolding which first removes a nonconstant mesoscopic intensity profile.

Let \(x_{1}\),\ldots,\(x_{d}\) be the eigenphases of a d-dimensional unitary U, measured in turns, so its eigenvalues are \(e^{2\pi ix_j}\). For a compactly supported continuous test \(\phi\) define the ordered factorial pair statistic \[
 C_d(U;\phi)=\frac1d\sum_{a\ne b}\sum_{n\in\mathbb Z}
              \phi\bigl(d(x_a-x_b+n)\bigr).                    \tag{FN.23}
\] Distinct indices, not distinct eigenvalues, are required in this definition. For \(L>0\) use the compact triangular test \[
 \phi_L(s)=\frac1L(1-|s|/L)_+,\quad
 \int\phi_L=1,\quad \phi_L(0)=1/L,\quad
 \widehat\phi_L(u)=\operatorname{sinc}^2(Lu)\geq0,               \tag{FN.24}
\] where \(\operatorname{sinc}(v)=\sin(\pi v)/(\pi v)\), with value 1 at zero. Fourier transforms use \(e^{-2\pi ius}\). Absolute convergence of the periodized Fourier series (the coefficients decay quadratically) gives the exact identity \[
 C_d(U;\phi_L)=d^{-2}\sum_{t\in\mathbb Z}
   \widehat\phi_L(t/d)(|\operatorname{Tr}U^t|^2-d).              \tag{FN.25}
\] If \(d>L\), the total diagonal subtraction is exactly 1/L, because \(\sum_n\phi_L(dn)=\phi_L(0)\). For any distinct positive integers \(r_{1}\),\ldots,\(r_{M}\), positivity therefore gives \[
 C_d(U;\phi_L)\geq1-\frac1L+
   2\sum_{j=1}^M\operatorname{sinc}^2(Lr_j/d)
                   |d^{-1}\operatorname{Tr}U^{r_j}|^2.         \tag{FN.26}
\] No control or truncation of the remaining high-frequency traces is needed for this lower bound; their contribution is nonnegative.

There is an equivalent counting interpretation. Let \(N_{L}\)(y) count the phases in the circular interval {[}y,y+L/d), with y uniform in {[}0,1). Then \(E_{y}\) \(N_{L}\)=L and an overlap-of-intervals calculation gives \[
 \operatorname{Var}_y N_L
      =L^2\bigl(C_d(U;\phi_L)-1\bigr)+L.                       \tag{FN.27}
\] This identity is also verified independently by integrating the piecewise-constant counting function, rather than reusing the pair sum.

\begin{theorem}[obstruction and necessary kick scale]\label{pass1:ch07_floquet_resonance-5} For the kicked Pell family of Theorem~\ref{pass1:ch07_floquet_resonance-1}, put \(\tau_K=\varepsilon_K\sqrt{2K}\). If \(\tau_{K}\) tends to a finite nonzero \(\tau\), then for every fixed \(L>0\), \[
 \liminf_{K\to\infty} C_{d_N}(U_{\varepsilon_K,N};\phi_L)
      \geq \Theta(\tau)-\frac1L,\qquad
 \Theta(\tau)=\sum_{m\in\mathbb Z}e^{-m^2\tau^2}>1,             \tag{FN.28}
\] and \[
 \liminf_{K\to\infty}\operatorname{Var}_y N_L
      \geq L^2\bigl(\Theta(\tau)-1\bigr).                      \tag{FN.29}
\] If \(\tau_K\to0\), then \(C_{d_N}(U_{\varepsilon_K,N};\phi_L)\to+\infty\) for every fixed \(L\). Consequently, convergence of the globally unfolded pair measures to a bounded density \(\rho(s)\) with \(\rho(s)\to1\) as \(|s|\to\infty\) requires \[
             |\varepsilon_K|\sqrt K\longrightarrow\infty.     \tag{FN.30}
\] This is a necessary condition only, not a sufficient universality criterion.
\end{theorem}

\textbf{Proof.} The recurrence trace limit remains valid when \(\tau_K\to\tau\) varies with \(K\). Lemma~\ref{pass1:ch07_floquet_resonance-3} bounds \(\tau_-(|X_K|)\), and \(|\tau_-(e^{im\tau_KX_K})-\tau_-(e^{im\tau X_K})|\leq |m|\,|\tau_K-\tau|\,\tau_-(|X_K|)\). The error in (FN.21) is uniform for bounded \(\tau_K\). Thus the normalized \(mK\)-trace tends to \(e^{-m^2\tau^2/2}\) for each fixed \(m\), including \(\tau=0\).

Apply (FN.26) with \(r_j=jK\). Since \(K/d_N\to0\), the sinc factors tend to one. At each fixed \(M\), the resulting lower bound is \(1-1/L+2\sum_{m=1}^M e^{-m^2\tau^2}\). Let \(M\) tend to infinity \emph{after} taking the lower limit in \(K\). This proves (FN.28), or divergence when \(\tau=0\). Equation (FN.27) gives (FN.29).

Suppose the pair measures converge vaguely to \(\rho\)(s)ds, where \(\rho\) is bounded and clusters to unit intensity. For every fixed L the limiting triangular statistic would be \[
 A_L=\int_{-1}^1(1-|u|)\rho(Lu)\,du.
\] Dominated convergence gives \(A_L\to1\) as \(L\to\infty\). But (FN.28) requires \(A_L\geq\Theta(\tau)-1/L\) for every \(L\), a contradiction. For \(\tau=0\), a finite locally bounded target already contradicts the fixed-\(L\) divergence. If \(|\varepsilon_K|\sqrt K\) does not tend to infinity, a bounded subsequence of \(\tau_K\) has a further subsequence converging to a finite limit. The same contradiction applies. End of proof.

The use of a compact test is essential for the logical scope. A Gaussian test without an independently controlled tail would only obstruct a Schwartz-test limit, not automatically vague local convergence. If the pair convergence definition uses only smooth compact tests, positive measures acquire locally uniform mass bounds by testing a bump which is one on a given compact set; approximation then extends convergence to the continuous triangular tests. Alternatively, start directly with a smooth nonnegative compact f of integral one, take \(\phi\)=f*tilde(f), and scale it: \(\widehat\phi=|\widehat f|^2\geq0\) and the same proof uses \(\phi\)(0)/L instead of 1/L. No unproved high-frequency-tail assumption is hidden in Theorem~\ref{pass1:ch07_floquet_resonance-5}.

The Poisson pair density \(\rho\)=1 and the sine-kernel density \(\rho_2(s)=1-\operatorname{sinc}^2(s)\) are both excluded. More explicitly, keeping only the first resonance gives \[
 \liminf C_{d_N}(U_{\varepsilon_K,N};\phi_L)
        \geq1+2e^{-\tau^2}-1/L.                               \tag{FN.31}
\] Since \(\int\phi_L\rho_2\leq1\), taking \(L=e^{\tau^2}\) yields a gap of at least \(e^{-\tau^2}\) from the sine-kernel prediction. At \(\tau=1\), already L=2 gives lower bound 1.2357588823\ldots, greater than both that prediction and the Poisson value 1.

The standard GOE bulk pair density is excluded by the bounded-clustering part of the theorem, not by an assumed pointwise bound \(\rho\leq1\). In unit intensity normalization it is the even continuation of \[
 \rho_1(s)=1-S(s)^2-S'(s)\int_s^\infty S(t)\,dt,\qquad
 S(s)=\operatorname{sinc}(s),\quad s\geq0.                       \tag{FN.32}
\] This standard formula is recorded in Brezin--Hikami, equation (181) after unfolding. It is bounded at zero, and integration by parts in the sine tail shows \(\rho_1(s)=1+O(s^{-2})\) at infinity. The same elementary criterion therefore applies. The random-matrix formulas, Fourier positivity and counting-variance identity are standard ingredients; no priority is claimed for those general tools.

This changes the route toward a positive BGS result: the critical and subcritical globally unfolded odd spectra considered here are not merely unproved universal families; they are \textbf{excluded} from the bounded clustering pair limits just specified. The statement does not settle nearest-neighbour distributions individually, density-adapted re-unfolding, a further desymmetrized block, or fixed nonzero classical kicks. It is not a refutation of a correctly qualified generic BGS conjecture.

The finite verifier keeps forty configurations (K=3,5,7, several \(\tau\) and L) and two synthetic controls. At K=7,\(\tau=1\),L=2 it numerically gives \(C_{d}\)=1.30936229449\ldots{} and Var \(N_{L}\)=3.23744917797\ldots, while the finite lower bound retaining K and 2K is 1.21861711382\ldots{} . These are floating-point checks of finite identities, not the proof of (FN.28). In particular the asymptotic bound is not incorrectly enforced on small K.

\part*{Chapter 8 --- Weighted echoes and Edgeworth corrections}
\addcontentsline{toc}{part}{Chapter 8 --- Weighted echoes and Edgeworth corrections}

\section{Weighted echoes, a resonance comb, and fixed-polynomial corrections}
\label{sec:weighted}
This section consolidates the weighted-observable, echo, resonance-comb,
and polynomial-correction notes. All statements use the fixed Pell map and
odd prime K, with $N+F\sqrt2=(1+\sqrt2)^K$, $D=(N-1)/2$,
$B_j=U_0^jBU_0^{-j}$, and $\tau_-=D^{-1}\operatorname{Tr}_-$.
The parameter $\tau$ below is a real kick amplitude, not a trace symbol.
The preceding normal form and lifting lemmas supply the stated inputs;
no independence of the finite Weyl operators is assumed.

\subsection{Bounded arrays and the ordered Gaussian law}
For real arrays f define
\[
 X_K(f)=(2K)^{-1/2}\sum_{j<K} f_j B_j,\qquad
 Q_K(f,g)=K^{-1}\sum_{j<K}f_jg_j.
\]
For fixed r and arrays uniformly bounded by M, the ordered mixed moments obey
\begin{equation}\label{eq:weighted-moments}
 \tau_-\bigl(X_K(f_1)\cdots X_K(f_{2r})\bigr)
 =\sum_{\pi\in\mathcal P_2(2r)}\prod_{\{a,b\}\in\pi}Q_K(f_a,f_b)
       +O_{r,M}(K^{-1})+O_{r,M}(K^r/N).
\end{equation}
For each fixed odd m, after the lifting threshold,
\[
 |\tau_-(X_K(f_1)\cdots X_K(f_m))|
       \leq M^m(2K)^{m/2}/(N-1).
\]
Thus convergence of a finite set of empirical Gram matrices gives the Wick
moments of a commuting real Gaussian vector with that covariance.

\begin{proof}
The exact parity-projected word trace is
$e_N(\phi)(N\mathbf1_{v=0}-1)/(N-1)$. It bounds the noncentral contribution
as displayed. The largest-gap norm lift and bounded-cluster argument confine
nonpaired words to $O_r(K^{r-1})$ choices. Odd central words are impossible:
a lifted polynomial is divisible by $z^2-6z+1$, so its value at one is
divisible by four and cannot be a sum of an odd number of signs.
For each matching of a balanced word, pair-label summation with the phase
replaced by one gives exactly the product of Q's. Repeated labels cost
$O_{r,M}(K^{-1})$. Its crossing phase is a product of $z_{a-b}^{\pm1}$.
Each individual difference is uniform, and $\sum_d|z_d-1|<9$; therefore
$|\prod u_e-1|\leq\sum_e|u_e-1|$ bounds its weighted mean error by
$9M^{2r}$ times the number of crossings divided by K. Distinct crossings
need not be independent. Summation proves the formula.
\end{proof}
The exact four-letter support theorem also gives
\begin{equation}\label{eq:weighted-comm}
 \|[X_K(f),X_K(g)]\|_{2,-}
 \leq\frac{6}{\sqrt K}\sqrt{\frac{2N}{N-1}}
                   \|f\|_\infty\|g\|_\infty,\qquad K\geq11.
\end{equation}
Indeed, the coefficient of $[B_t,B_s]$ is
$(f_tg_s-f_sg_t)/(2K)$, bounded by $\|f\|_\infty\|g\|_\infty/K$.
The four Weyl terms for different unordered pairs have disjoint support.
Their full-space squared norm is therefore at most
$4\|f\|_\infty^2\|g\|_\infty^2K^{-2}\sum_{j=1}^{K-1}(K-j)a_j^2$.
Use $a_{K-j}=a_j$, $a_j^2\leq2a_j$, and $\sum a_j<9$, then restrict the
positive quadratic form to the odd sector. This proves
\eqref{eq:weighted-comm}; it is not an operator-norm assertion.

\subsection{The interaction-picture echo}
Let $U_\varepsilon=e^{i\varepsilon B}U_0$,
$\varepsilon=\tau/\sqrt{2K}$ and $E_K(t)=U_\varepsilon^tU_0^{-t}$.
Write $t=mK+s$, $0\leq s<K$, including negative t by the usual floor
convention. Define $w_{t,K}(j)=m+\mathbf1_{j<s}$.
The centered coefficients $c_j=(K-1-2j)/(2K)$ satisfy the exact identity
\[
 c_j+t/K-c_{(j-t)\bmod K}=w_{t,K}(j).
\]
Inserting the finite normal form, telescoping t unitary factors, and combining
the three exponentials with \eqref{eq:weighted-comm} gives
\begin{equation}\label{eq:echo-finite}
 \|E_K(t)-e^{i\tau X_K(w_{t,K})}\|_{2,-}
       \leq (11C+2)\tau^2/\sqrt K\qquad(|t|\leq CK).
\end{equation}
To detail the constant, the original normal-form error contributes
$6|t/K|\tau^2/\sqrt K$. The comparison has three factors with arrays c,
the constant $t/K$, and the negative shifted c. The two Duhamel errors sum
to at most $c_N(3|t/K|+3/4)\tau^2/\sqrt K$, where
$c_N=\sqrt{2N/(N-1)}\leq\sqrt{7/3}$. This yields the displayed bound.
The same argument uses adjoints at negative times.

Let $\mathcal B(a)$, $0\leq a\leq1$, be one standard Brownian motion and,
for $\alpha=m+a$ with $0\leq a<1$, put
\[
 Y(\alpha)=m\mathcal B(1)+\mathcal B(a).
\]
For $\beta=n+b$ its covariance is
$mn+mb+na+\min(a,b)$. The arrays $w_{t,K}$ have precisely this limiting
Gram matrix when $t/K\to\alpha$. Therefore for every fixed finite collection
of integer exponents $r_l$ and times $t_l/K\to\alpha_l$,
\[
 \tau_-\left(\prod_l E_K(t_l)^{r_l}\right)
 \longrightarrow
 \exp\left[-\frac{\tau^2}{2}\sum_{l,j}r_lr_j
                     \operatorname{Cov}(Y(\alpha_l),Y(\alpha_j))\right].
\]
The order of the factors on the left is retained. Commutator estimates allow
them to be combined in the limit; there is no common eigenbasis at finite K.
In particular the phase variance of one echo is
$\tau^2[m^2+(2m+1)a]$. The same Brownian interval is reused on each cycle:
fresh independent noise would give a different variance. This is a tracial
finite-dimensional-distribution statement, not convergence of quantum
measurement trajectories in a path space.

\subsection{The full bounded recurrence window}
For $1\leq s<K$, the fixed-point equations for $W(v)U_0^s$ and
$RW(v)U_0^s$ have unique solutions, so
$|\tau_-(W(v)U_0^s)|\leq1/D$. Hence, with $X_K=X_K(1)$,
\[
 |\tau_-(X_K^lU_0^s)|\leq(2K)^{l/2}/D.
\]
For $|t|\leq CK$, $K\nmid t$, let $A_0=C|\tau|$. Taylor's integral
remainder of degree $2r-1$ and the normal form yield
\[
 |\tau_-(U_\varepsilon^t)|
 \leq\frac{6C\tau^2}{\sqrt K}
 +\frac1D\sum_{l<2r}\frac{(A_0\sqrt{2K})^l}{l!}
 +\frac{A_0^{2r}}{(2r)!}\tau_-(X_K^{2r}).
\]
First let K grow at fixed r, then r grow. The last limiting bound is
$A_0^{2r}/(2^rr!)\to0$. Thus nonresonant traces vanish uniformly in this
window, whereas $\tau_-(U_\varepsilon^{mK})\to e^{-m^2\tau^2/2}$.
If $e^{i\theta_l}$ are the actual eigenvalues, the empirical measure of
$(e^{i\theta_l},e^{iK\theta_l})$ converges to Haar measure times the wrapped
normal law: its Fourier coefficients are traces at $a+bK$.
This is a two-scale law of one eigenvalue, not independence of two levels.

\subsection{Minimal Pell relations and the first polynomial correction}
For an integer polynomial P, write $\|P\|_1$ for its coefficient norm.
If $P\ne0$ and $P(\lambda)=0$, then $\|P\|_1\geq8$, with equality only for
$\pm z^j(z^2-6z+1)$. Indeed $P(1)$ is divisible by four, so its norm is even;
if it were at most six the leading term would dominate since $\lambda>5$.
At norm eight the leading coefficient is one in absolute value. After a sign
change write $P=z^L+c z^{L-1}+R$, $\|R\|_1=7-|c|$. Evaluation gives
$|(6+c)\lambda-1|\leq7-|c|$, forcing $c=-6$ and then
$R=z^{L-2}$. Thus the first genuinely nonpaired central words have length eight.
The bounded-cluster argument improves their normalized contribution at any
fixed even order to $O_r(K^{-3})$.

For a fixed simple crossing graph, expand the phase product through first
order in $z_d-1$. Any two distinct edges form a forest, so their label
differences are independent uniform variables. The remainder is therefore
$O(K^{-2})$. A matching with c crossings contributes
$1+c(S_K-1)+O(K^{-2})$, $S_K=K^{-1}\sum_dz_d$. The total crossing count is
$A_r=\binom{2r}{4}(2r-5)!!$: choose four interlaced endpoints and match the
rest. Four repeated positions cause one excess matching; their six balanced
sign assignments contribute $-3A_r/(2K)+O(K^{-2})$ to the moment: the
matching-indexed sum counts each such word with one excess matching, so
the excess is subtracted. The crossings contribute
$A_r(S_K-1)=-(\gamma-3/2)A_r/K+O(K^{-2})$, and the two corrections add
to $-\gamma A_r/K$. Since
$S_K=1-(\gamma-3/2)/K+O((KN)^{-1})$, it follows that
\begin{equation}\label{eq:edgeworth-mom}
 \tau_-(X_K^{2r})=(2r-1)!!-
       \frac{\gamma}{K}\binom{2r}{4}(2r-5)!!
       +O_r(K^{-2})+O_r(K^r/N).
\end{equation}
Here
\[
 \gamma=\frac32+2\sum_{j\geq0}
 \left[1-\cos\left(\frac{4\pi}{\sqrt2(1+\sqrt2)}\lambda^{-j}\right)\right].
\]
For every fixed real polynomial f and standard Gaussian Z,
\[
 \tau_-(f(X_K))=\mathbb Ef(Z)-\frac{\gamma}{24K}\mathbb Ef^{(4)}(Z)
             +O_f(K^{-2})+O_f(K^{\deg f/2}/N).
\]
This is a fixed-polynomial expansion, not a distribution-function Edgeworth
theorem. Constants are not controlled as the polynomial degree grows.

For completeness, here is the stronger crossing-graph expansion, with the
notation it needs. For a fixed simple graph $G$ on the pair labels, let
the labels $t_v$ be independent and uniform in $\mathbf Z/K\mathbf Z$,
orient each edge, choose signs $\sigma_e\in\{-1,1\}$, and put
$Z_G=\prod_{e=(u,v)}z_{t_v-t_u}^{\sigma_e}$, where $z_d$ are the crossing
phases above; $\mathbb E$ is the average over the labels. Then
\[
 \mathbb E Z_G=S_K^{|E(G)|}+O_G(K^{-3}+N^{-1}).
\]
Indeed, expand $\prod_e(1+(z_e^{\sigma_e}-1))$. Every subset of at most
two edges is a forest, so its edge differences are independent and
uniform, and its expectation is exactly $(S_K-1)$ per edge, resp.\
$(S_K-1)^2$ per pair; this matches the binomial expansion of
$S_K^{|E(G)|}$ through second order, the remaining binomial terms being
$O_G(K^{-3})$ since $S_K-1=O(K^{-1})$. A subset of at least three edges
either contains a three-edge forest, whose expected absolute defect
product is at most $(9/K)^3$ by $\sum_d|z_d-1|<9$, or is a triangle. For
a triangle the bound is on the defect product, not on the unit-modulus
phases themselves: the three differences sum to zero modulo $K$, their
signed representatives therefore sum to $0$ or $\pm K$ (for $\pm K$ all
three have the same sign, since each is at most $(K-1)/2$ in absolute
value), in both cases
$\sum_i|\delta_i|\le K$, i.e.\ $\sum\ell\geq(K-3)/2$ for
$\ell(d)=(K-1)/2-|\delta|$; the phase-rate bound
$|z_d-1|\le A\lambda^{-\ell(d)}/(1-\lambda^{-K})$ with
$A=4\pi/(\sqrt2(1+\sqrt2))$ then gives a defect product
$O(\lambda^{-K/2})=O(N^{-1})$, uniformly in the labels, and if one
difference is zero the product is zero. All remaining defect factors are
bounded by two. This proves the stated refinement.

At order eight, the minimal relation is
$(3,-4)+(3,4)-6(1,0)=0$. Among its 56 orders the seven half-symplectic phases
$-12,-8,-4,0,4,8,12$ each occur eight times. Both signs and K time shifts
therefore give the isolated nonpaired central contribution
\[
 \frac{N}{N-1}K^{-3}
  [1+2\cos(8\pi/N)+2\cos(16\pi/N)+2\cos(24\pi/N)].
\]
It is $7/K^3+O((K^3N)^{-1})$, in addition to all balanced contributions;
it is not the whole eighth moment. The sufficient eight-letter lifting
threshold holds for odd primes $K\geq29$. Its remaining odd-parity trace
error is at most $16K^4/(N-1)$.

\subsection{Subcritical microscopic concentration without a Gaussian input}
For a circle spectrum define $a^2=D^{-1}\sum_l\operatorname{dist}
(\theta_l,K^{-1}\mathbb Z)^2$. Arcs of halfwidth $w<1/(2K)$ about the roots
contain at least $D(1-a^2/w^2)_+$ points. Counting gaps and applying
Cauchy--Schwarz to cells of length $x/D$ gives
\[
 \frac1D\#\{S_l>x\}\leq\frac{2a^2}{w^2}+\frac KD+\frac{2Kw}{x},
 \qquad
 R_D(x)\geq\frac{(1-a^2/w^2)_+^2}{K/D+2Kw/x}-1.
\]
An eigenvalue matching for normal matrices follows by decomposing their
squared-eigenbasis-overlap matrix into permutations. Chord distance bounds
circle distance, and the finite normal form gives
\[
 a\leq\frac{3\varepsilon^2+
            (|\varepsilon|/\pi)\sqrt{N/(N-1)}}{\sqrt K}.
\]
Consequently $|\varepsilon|\sqrt K\to0$ implies $Ka\to0$, collapse of the
gap law to $\delta_0$, and divergence of $R_D(x)$ for each fixed x. Set
$w=(Ka)^{2/3}/K$ for the gap estimate and $w=2a$ for the pair estimate;
the case $a=0$ is direct. These conclusions need no Gaussian limit and
refer to the original phases, not folded phases or isolated old eigenspaces.

\part*{Chapter 7 --- Floquet, coboundary and smooth resonance fields (continued)}

\section{A second-order recurrence scale for a coboundary kick}\label{sec:coboundary}
Set \[
 T=\begin{pmatrix}3&2\\4&3\end{pmatrix},\qquad
 \alpha=1+\sqrt2,\quad \lambda=3+2\sqrt2,\quad
 N+F\sqrt2=\alpha^K,
\] where K tends to infinity through odd primes. Then N is odd, \(N^2-2F^2=-1\), and \(N\asymp\lambda^{K/2}\). The moduli need not be prime. Use the phase-normalized cat quantization \(\widetilde U_0\) from the \hyperref[sec:floquet]{finite Floquet construction}: \(\widetilde U_0\) is the operator denoted \(U_{N,\mathrm{Fl}}=VU_0V^*\) there, not the auxiliary \(U_0=S_h^{-1}\) of (FN.1); the tilde distinguishes the two throughout this coboundary section. It satisfies \(\widetilde U_0W(v)\widetilde U_0^*=W(Tv)\) and \(\widetilde U_0^K=I\). Let \[
 W(a,b)f(x)=e_N(a(x-b/2))f(x-b),\qquad B_v=W(v)+W(-v).
\] All symmetrized observables and \(\widetilde U_0\) preserve odd parity. Write \(d=(N-1)/2\) and \(\tau_-(M)=d^{-1}\operatorname{Tr}_{-}M\).

The perturbation is now a different, fixed Fourier observable: \[
 A=B_{(1,0)},\quad D=\widetilde U_0A\widetilde U_0^*=B_{(3,4)},\quad
 H=A-D,\quad V_\varepsilon=e^{i\varepsilon H}\widetilde U_0.                 \tag{CB.1}
\] It is a coboundary: \(\sum_{j=0}^{K-1}\widetilde U_0^jH\widetilde U_0^{-j}=0\) exactly. The real classical Fourier symbol of H is independent of N. The products of the two coordinates of (1,0), (3,4), (2,4), and (4,4) are even, so the centered finite-Weyl convention introduces no N-dependent sign for these fixed Fourier modes. As throughout the project, the quantum exponent \(\varepsilon\) corresponds to a classical action amplitude proportional to \(\varepsilon\)/N, not \(\varepsilon\) itself.

Put \[
 s_N=\sin(4\pi/N)>0,\quad \zeta=\varepsilon^2s_N,\quad
 t_K=2\zeta\sqrt K.
\]
\begin{theorem}[second-order critical law]\label{pass1:ch07_coboundary_resonance-1} If \(t_K\to t\in[0,\infty)\), then for every fixed integer m, \[
 \tau_-(V_{\varepsilon_K}^{mK})\longrightarrow e^{-m^2t^2/2}.     \tag{CB.2}
\] Thus the empirical spectrum of \(V_{\varepsilon_K}^K\) tends to the wrapped normal law of variance \(t^2\), or to the point mass at 1 when t=0. The critical quantum scale is \[
 |\varepsilon_K|=\left(\frac{t}{2s_N\sqrt K}\right)^{1/2}
 \sim\sqrt{\frac{t}{8\pi}}\,\sqrt N K^{-1/4}\quad(t>0).         \tag{CB.3}
\] This can grow without bound, even though the classical amplitude tends to zero. The sign of \(\varepsilon\) does not change the limit in (CB.2).
\end{theorem}

\begin{theorem}[negative microscopic consequences]\label{pass1:ch07_coboundary_resonance-2} For the original propagator in (CB.1), on the whole odd space and with global mean-spacing normalization, the following hold.
\end{theorem}

If \(|\varepsilon_K|K^{1/4}/\sqrt N\to0\), the empirical nearest-neighbour spacing measures tend weakly to \(\delta_0\), and the normalized number of ordered distinct pairs at distance at most x mean spacings tends to infinity for every fixed \(x>0\).

For any subsequence on which \(|\varepsilon_K|K^{1/4}/\sqrt N\) is bounded, the pair measures cannot converge to a bounded density R(s) with \(R(s)\to1\) as \(|s|\to\infty\). In particular, escape from this obstruction requires \[
 |\varepsilon_K|K^{1/4}/\sqrt N\longrightarrow\infty.            \tag{CB.4}
\] Condition (CB.4) is only necessary. No positive universality theorem beyond it is proved. These statements concern global unfolding of the specified odd spectrum, not arbitrary clusterwise unfolding or an assertion that all arithmetic symmetries are resolved. They do not refute fixed-system BGS.

\subsection{A second-order loop estimate valid at large quantum exponents}\label{coboundary:a-second-order-loop-estimate-valid-at-large-quantum-exponents}

\begin{lemma}\label{pass1:ch07_coboundary_resonance-3} For finite Hermitian A,D and any real \(\varepsilon\), \[
 \left\|e^{-i\varepsilon A}e^{i\varepsilon(A-D)}e^{i\varepsilon D}
       -e^{-\varepsilon^2[A,D]/2}\right\|_{\rm op}
 \leq |\varepsilon|^3\left(
       \frac12\|[A,[A,D]]\|_{\rm op}
       +\frac16\|[D,[A,D]]\|_{\rm op}\right).                  \tag{CB.5}
\] There is no hypothesis \(|\varepsilon|\|A\|\ll1\).
\end{lemma}

\textbf{Proof.} Write \(Q=[A,D]\) and \(L(u)=e^{-iuA}e^{iu(A-D)}e^{iuD}\). Direct differentiation gives \[
 L'(u)L(u)^*=-\int_0^u
       \operatorname{Ad}_{e^{-iuA}e^{iv(A-D)}}Q\,dv.
\] The logarithmic derivative of \(E(u)=e^{-u^2Q/2}\) is -uQ. By differentiating unitary conjugation and integrating, \[
 \|\operatorname{Ad}_{e^{-iuA}e^{iv(A-D)}}Q-Q\|
 \leq u\|[A,Q]\|+v\|[A-D,Q]\|\quad(0\leq v\leq u).
\] Duhamel comparison of the two unitary evolutions, integrated over \(0\leq v\leq u\leq\varepsilon\), is at most \(\varepsilon^3(\|[A,Q]\|/3+\|[A-D,Q]\|/6)\). The triangle inequality proves (CB.5) for nonnegative \(\varepsilon\). Replace A,D by -A,-D for the other sign. This argument is a norm estimate, not a truncation of a formal BCH series. End of proof.

Let u=(2,4), z=(4,4) and C=\(B_{u}\)-\(B_{z}\). Since \(W(v)W(w)=e_N(\omega(v,w)/2)W(v+w)\) and \(\omega((1,0),(3,4))=4\), exact multiplication yields \[
 [A,D]=-2i s_N C.                                             \tag{CB.6}
\] Each double commutator in (CB.5) has at most eight Weyl terms of coefficient modulus \(4s_N^2\). Therefore both norms are at most \(32s_N^2\), and \[
 \left\|e^{-i\varepsilon A}V_\varepsilon e^{i\varepsilon A}
                  -e^{i\zeta C}\widetilde U_0\right\|_{\rm op}
 \leq R_N(\varepsilon):=\frac{64}{3}|\varepsilon|^3s_N^2.        \tag{CB.7}
\] The same bound holds in normalized odd Hilbert--Schmidt norm. At the scale (CB.3), \(KR_N=O_t(K^{1/4}/\sqrt N)\); this tends to zero. The small quantity is controlled by nested commutators, not by the norm of \(\varepsilon\) A.

\subsection{A two-species arithmetic lemma: collisions must not be discarded}\label{coboundary:a-two-species-arithmetic-lemma-collisions-must-not-be-discarded}

Identify an integer vector v=(a,b) with \(\psi(v)=b+\sqrt2 a\). Then \[
 \psi(Tv)=\lambda\psi(v),\quad
 \psi(u)=4+2\sqrt2,\quad \psi(z)=\sqrt2\psi(u),\quad
 \operatorname{Norm}(\psi(u))=8.
\] Since N is odd, multiplication by \(\psi(u)\) is invertible in \((\mathbb Z/N\mathbb Z)[\sqrt2]\). The two species of time translates are therefore faithfully represented by \(\lambda^j\) and \(\sqrt2\lambda^j\). Also \(T^K=I\) modulo N, as follows from \(Q_{2K}=2N^2+1\) and \(P_{2K}=2NF\). Rotating all labels modulo K preserves a vector relation.

\begin{lemma}[lifting]\label{pass1:ch07_coboundary_resonance-4} A signed m-atom modular zero relation in these two species, after cancellation of identical opposite atoms and rotation across a largest cyclic gap, has the form \[
 P(\lambda)=\sum_\nu \epsilon_\nu b_\nu\lambda^{j_\nu},
 \quad b_\nu\in\{1,\sqrt2\},\quad
 0\leq j_\nu\leq K-\lceil K/m\rceil.
\] For sufficiently large K at fixed m this expression is zero in \(\mathbb Z[\sqrt2]\), not just modulo N. For m=4, \(K\geq11\) suffices.
\end{lemma}

\textbf{Proof.} There are at most m occupied time labels, hence a cyclic gap of at least ceil(K/m). Modular vanishing and the invertibility just proved imply \(P(\lambda)\in N\mathbb Z[\sqrt2]\). Its nonzero algebraic norm would thus have absolute value at least \(N^2\), whereas the two real embeddings give \[
 |\operatorname{Norm}P(\lambda)|
 \leq2m^2\lambda^{K-\lceil K/m\rceil}.
\] Since \(N^2=\lambda^K(1-\lambda^{-K})^2/4\), the latter is smaller than \(N^2\) for large K at fixed m. For m=4 it suffices that \(\lambda^{\lceil K/4\rceil}(1-\lambda^{-K})^2>128\). At \(K\geq11\) this follows from \(\lambda>11/2\), \(\lceil K/4\rceil\geq3\), and \(1-\lambda^{-K}>9/10\), since \((11/2)^3(9/10)^2>128\). End of proof.

A four-atom relation in characteristic zero need \textbf{not} balance within individual species. Apart from cancellation, the sole type is \[
 z_{j+1}-u_{j+1}=z_j+u_j,\qquad v_j:=T^jv,                     \tag{CB.8}
\] up to translation and multiplication by -1. The identity is already true over the integers, since \(T(z-u)=z+u\).

Here is a direct completeness proof. Look at the maximal exponent E after cancellation. With one atom at E, its coefficient has modulus at least 1, whereas the three lower atoms sum to at most \(3\sqrt2/\lambda<1\) times \(\lambda^E\). With three atoms at E, their nonzero coefficient has modulus at least \(2-\sqrt2>\sqrt2/\lambda\), which the remaining lower atom cannot cancel. With four atoms at E, linear independence of 1 and \(\sqrt2\) forces balance. With two atoms at E, the only possible nonzero coefficient small enough is \(\pm(\sqrt2-1)\). The lower atoms must both be at E-1: otherwise at that scale their modulus is at most \(\sqrt2(1+\lambda^{-1})<1+\sqrt2\). At E-1 they must be one of each species with equal signs because \(\lambda(\sqrt2-1)=1+\sqrt2\). This is exactly (CB.8). Relations with fewer occurrences reduce to the same cases; there is no nontrivial relation with at most three atoms.

\begin{corollary}[multiplicity two]\label{pass1:ch07_coboundary_resonance-5} For \(K\geq11\), form all signed unordered pairs with distinct time labels: \[
 \epsilon v_t+\delta w_s,
 \quad 0\leq t<s<K,\quad v,w\in\{u,z\},\quad
 \epsilon,\delta\in\{-1,1\}.
\] There are \(8K(K-1)\) such pairs. Their vector keys are nonzero, have multiplicity at most two, and have precisely \[
 4K\ \hbox{double keys},\quad 8K(K-2)\ \hbox{single keys},
 \quad 8K^2-12K\ \hbox{distinct keys}.                          \tag{CB.9}
\] The double classes are the two identities \[
 z_{j+1}-z_j=u_{j+1}+u_j,\qquad
 z_{j+1}-u_j=u_{j+1}+z_j,
\] with their negatives, for j modulo K. The third partition of (CB.8) uses equal-time pairs and is not in the domain. Each displayed signed pair belongs to only one class: its species, the two adjacent cyclic labels, and its signs determine the other pair. Thus there are no triple classes or overlapping equivalence chains. Lemma~\ref{pass1:ch07_coboundary_resonance-4} ensures that no additional modular relation is missed. K=3 is a useful negative control: it has six triple keys, and the range restriction must not be removed.
\end{corollary}

\subsection{Weighted phase control and a finite normal form for C}\label{coboundary:weighted-phase-control-and-a-finite-normal-form-for-c}

Write \(C_j=\widetilde U_0^jC\widetilde U_0^{-j}\), \(\bar C=K^{-1}\sum_j C_j\), \(G_C=\sum_j c_jC_j\), where \(c_j=(K-1-2j)/(2K)\). Then \[
 G_C-\widetilde U_0G_C\widetilde U_0^*=C-\bar C,\qquad [\bar C,\widetilde U_0]=0.                \tag{CB.10}
\] Define a phase majorant, including both species, \[
 a_j=\max_{v,w\in\{u,z\}}
          |e_N(\omega(v,T^jw))-1|.
\] One has \(a_{K-j}=a_j\). In particular, no independence assumption about the several phase factors is being made.

For a centered representative \(\delta\) of j modulo K, \[
 \omega(v,T^\delta w)
 =(2v_1w_1-v_2w_2)P_{2\delta}+\det(v,w)Q_{2\delta}.
\] For the four choices of species these are \(-8P_{2\delta}\), \(16P_{2\delta}\), and \(\pm8Q_{2\delta}\), respectively. Their absolute values are at most \(8\lambda^{|\delta|}\). If \(|\delta|=(K-1)/2-\ell>0\), this gives \[
 a_j\leq\frac{32\pi}{\alpha(1-\lambda^{-K})}\lambda^{-\ell}.
\] At zero, \(a_0\leq16\pi/N\). Summing both signs, including their weights \(\ell/K\) and \((\ell+1)/K\), proves the explicit estimates \[
 \begin{split}
 L_K:=\sum_j a_j
   &\leq\frac{32\pi}{1-\lambda^{-K}}+\frac{16\pi}{N}<128,\\
 K\sum_j|c_j|a_j
   &\leq\frac{16\sqrt2\pi}{1-\lambda^{-K}}+
                                      \frac{8\pi K}{N}<128.
 \end{split}                                                    \tag{CB.11}
\] For \(K\geq3\), use \(\lambda^{-K}<1/10\), \(N\geq2K+1\), \(\pi<22/7\), and \(\sqrt2<3/2\) to check the last inequalities. These constants are deliberately conservative.

The commutator \([C_t,C_s]\) has sixteen Weyl terms, each of modulus at most \(a_{s-t}\). Corollary~\ref{pass1:ch07_coboundary_resonance-5} says that each key in the sum over \(t<s\) has at most two preimages. Weyl matrices are orthonormal in the \emph{full} normalized Hilbert--Schmidt inner product, also over composite odd N. Cauchy--Schwarz inside each key class therefore gives \[
 \begin{split}
 \|[G_C,\bar C]\|_{2,N}^2
 &\leq\frac{32}{K^4}\sum_{j=1}^{K-1}(K-j)j^2a_j^2\\
 &=\frac{16}{K^3}\sum_{j=1}^{K-1}j(K-j)a_j^2
 \leq\frac{1024}{K}.
 \end{split}                                                    \tag{CB.12}
\] The equality between the two sums uses \(a_{K-j}=a_j\); the final bound uses \(a_j^2\leq2a_j\) and (CB.11). Norm restriction to the odd subspace obeys \(\|M\|_{2,-}\leq\sqrt{2N/(N-1)}\|M\|_{2,N}\) by positivity. We have not assumed orthogonality of restricted Weyl matrices.

The triangle inequality in operator norm also gives \[
 \|[G_C,C]\|_{\rm op}\leq2048/K,\qquad
 \|[C,\bar C]\|_{\rm op}\leq2048/K.                           \tag{CB.13}
\] The usual two-exponential Duhamel inequality is \(\|e^{ixA}e^{ixD}-e^{ix(A+D)}\|\leq x^2\|[A,D]\|/2\) in any unitarily invariant norm. Applying it first to -\(G_{C}\),C, then to C-\(G_{C}\) and \(\widetilde U_0G_C\widetilde U_0^*\), using (CB.10)--(CB.13), proves \[
 \left\|e^{-i\zeta G_C}(e^{i\zeta C}\widetilde U_0)e^{i\zeta G_C}
                        -e^{i\zeta\bar C}\widetilde U_0\right\|_{2,-}
 \leq\frac{1000\zeta^2}{\sqrt K}\quad(K\geq11).               \tag{CB.14}
\] Indeed the less rounded bound is \(\zeta^2(2048/K+16\sqrt{2N/(N-1)}/\sqrt K)\), which is smaller than (CB.14) already using \(K\geq11\) and \(N\geq7\). This is another estimate valid for every real zeta, not a formal expansion requiring a small generator.

Combining the two gauges and telescoping unitary powers gives \[
 \left|\tau_-(V_\varepsilon^{mK})-
          \tau_-(e^{im\zeta K\bar C})\right|
 \leq |m|KR_N(\varepsilon)+1000|m|\zeta^2\sqrt K.              \tag{CB.15}
\]

\subsection{The two-species Gaussian limit}\label{coboundary:the-two-species-gaussian-limit}

\begin{lemma}\label{pass1:ch07_coboundary_resonance-6} Put \(X_K=(\sqrt K/2)\bar C\). For every fixed \(r\geq1\), \[
 \tau_-(X_K^{2r})=(2r-1)!!+O_r(K^{-1})+O_r(K^r/N).             \tag{CB.16}
\] For every fixed odd m, eventually \(|\tau_-(X_K^m)|\leq(4K)^{m/2}/(N-1)\). Consequently its empirical spectral measures converge weakly and in quadratic Wasserstein distance to standard normal.
\end{lemma}

\textbf{Proof.} There are 4K signed Weyl atoms, with coefficient +1 for species u and -1 for species z. For a word with central phase \(e_{N}\)(\(\phi\)) and total index v, the exact odd trace is \[
 e_N(\phi)\frac{N\mathbf1_{v=0}-1}{N-1}.                       \tag{CB.17}
\] This follows from \(\operatorname{Tr}W(v)=N\mathbf1_{v=0}\) and \(\operatorname{Tr}(RW(v))=1\). The noncentral contribution after normalization \((4K)^{-m/2}\) is at most \((4K)^{m/2}/(N-1)\).

Apply Lemma~\ref{pass1:ch07_coboundary_resonance-4} to central words. The following bounded-cluster argument handles the unavoidable identities (CB.8), rather than incorrectly excluding them. In a nonzero polynomial relation with at most m signed atoms and coefficients 1 or \(\sqrt2\), a gap d with \(\lambda^d>2m^2\) separates independently vanishing components. To see this, write \(P=P_{\rm low}+Z^{a+d}P_{\rm high}\) with degree of \(P_{low}\) at most a. If \(P_{\rm high}(\lambda)\ne0\), its nonzero integer norm and the bound \(|P_{\rm high}(\lambda^{-1})^{\rm conjugate}|\leq\sqrt2m\) give \(|P_{\rm high}(\lambda)|\geq1/(\sqrt2m)\). This contradicts cancellation by \(|P_{\rm low}(\lambda)|\leq\sqrt2m\lambda^a\). Iterating produces components with spans bounded solely in m.

Cancel opposite identical atoms first and restore these pairs after splitting. Every nonzero component then contains at least three occurrences. A two-atom relation cannot change species, because its two algebraic norms would be 1 and -2, and cannot change exponent, because \(\lambda>1\). A nonbalanced 2r-word thus has at most r-1 components. There are only finitely many partitions of its positions, species, signs and bounded offset patterns at fixed r. Each component has K possible cyclic translations, so there are \(O_r(K^{r-1})\) such words. Choosing the largest-gap cut introduces no additional K factor: this count is over the actual translated components. Their normalized contribution is \(O_r(1/K)\).

In a balanced word, pair identical species-time labels with opposite signs. For each pairing there are \(K^r\) time choices, \(2^r\) species choices and \(2^r\) sign choices. Their product is precisely the normalization \((4K)^r\). Repeated pair labels cause at most \(O_r(K^{r-1})\) overcount. Signs from species z occur twice in each pair and hence multiply to +1. The Weyl phase is a product of crossing factors or their conjugates. For fixed species, each crossing time difference is uniform on the cyclic group. By (CB.11) and \(|\prod u_j-1|\leq\sum|u_j-1|\), its average error is at most 128 times the number of crossings divided by K. No independence of the crossing factors is assumed. Summing over the \((2r-1)!!\) pairings gives (CB.16).

For odd m, reduce a lifted relation in \(\mathbb Z[\sqrt2]\) modulo 2. Here \(\lambda=3+2\sqrt2\) reduces to 1. The coefficients of both 1 and \(\sqrt2\) in the relation must be even. Both species counts would therefore be even, contradicting odd m. Thus no odd central word exists after lifting.

Bounded second moments yield tightness, higher even moments yield uniform integrability, and the Gaussian moment sequence is determinate, since its Carleman sum diverges. Therefore all weak subsequential limits are normal. Convergence of the second moments gives the claimed Wasserstein upgrade. End of proof.

Since \(\zeta K\bar C=t_KX_K\), (CB.15) and Lemma~\ref{pass1:ch07_coboundary_resonance-6} prove Theorem~\ref{pass1:ch07_coboundary_resonance-1}. More explicitly, the finite comparison error is bounded by \[
 |m|\left(KR_N(\varepsilon)+\frac{250t_K^2}{\sqrt K}\right).    \tag{CB.18}
\] For bounded \(t_{K}\), \(KR_N=O(K^{1/4}/\sqrt N)\) follows directly from \(\varepsilon^2=t_K/(2s_N\sqrt K)\) and \(s_N\leq4\pi/N\). The remaining Gaussian characteristic-function error is \(o(1)\); no universal \(K^{-1/2}\) rate is asserted for that separate step.

\subsection{Microscopic consequences without conflating unfolding conventions}\label{coboundary:microscopic-consequences-without-conflating-unfolding-conventions}

For completeness we give the two spectral implications used in Theorem~\ref{pass1:ch07_coboundary_resonance-2}. First, the 4K Weyl atoms in \(C_{j}\) are distinct, so \(\operatorname{Tr}(\bar C^2)/N=4/K\). Distinctness follows within species from the unit conditions on \(1\pm h^j\) in the Pell construction. Between species, norm 1 versus -2 would require N to divide 3, which is impossible for \(K\geq3\). Consequently \[
 \tau_-(\bar C^2)\leq\frac{8N}{K(N-1)}.
\] The squared-overlap matrix between orthonormal eigenbases of two unitaries is doubly stochastic. Decomposing it into permutations yields an eigenvalue matching whose mean squared chord cost is bounded by normalized squared Hilbert--Schmidt distance. Circle distance in \(\mathbb R/\mathbb Z\) is at most one quarter of chord distance. Using this matching, (CB.7),(CB.14), and \([\widetilde U_0,\bar C]=0\), the RMS distance a of the spectrum of \(V_{\varepsilon}\) from the K-point lattice satisfies \[
 a\leq\frac{R_N(\varepsilon)}4+
          \frac{250\zeta^2}{\sqrt K}+
          \frac{|\zeta|}{2\pi}\sqrt{\frac{8N}{K(N-1)}}.         \tag{CB.19}
\] When \(t_{K}\) tends to zero, Ka tends to zero, and K/d tends to zero.

For an arbitrary d-point circle multiset with this RMS lattice distance, placing arcs of half-width \(w<1/(2K)\) around the lattice and using Markov's inequality gives, for its unfolded cyclic gaps \(S_{l}\), \[
 d^{-1}\#\{l:S_l>x\}\leq2a^2/w^2+K/d+2Kw/x.                  \tag{CB.20}
\] There are at most \(da^2/w^2\) bad points; at most twice as many gaps touch one, at most K other gaps leave an arc, and the within-arc gaps have total length at most 2Kw. This proves (CB.20). Partitioning the arcs into cells of length at most x/d and applying Cauchy--Schwarz to their occupancies also gives the ordered close-pair bound \[
 R_2(x)\geq\frac{\max(0,1-a^2/w^2)^2}{K/d+2Kw/x}-1.           \tag{CB.21}
\] Take \(w=(Ka)^{2/3}/K\) for (CB.20), and w=2a for (CB.21). These prove subcritical spacing collapse and divergent close-pair counts. The a=0 case follows by counting at most K nonzero gaps and using \(R_2(x)\geq d/K-1\). The average spacing remains exactly one: weak convergence to \(\delta_{0}\) does not imply convergence of first moments, which are carried by rare long gaps.

For bounded nonzero critical parameter a different argument is necessary. Let \(\phi_L(s)=L^{-1}(1-|s|/L)_+\), a compactly supported nonnegative function with integral one and Fourier transform \(\operatorname{sinc}^2(L\xi)\). For the eigenangles \(\theta_{l}\) in \(\mathbb R/\mathbb Z\), define \[
 C_d(\phi_L)=\frac1d\sum_{l\ne j}\sum_{q\in\mathbb Z}
                       \phi_L(d(\theta_l-\theta_j+q)).
\] Poisson summation, including the zero Fourier mode and the removed diagonal, proves, for \(d>L\), \[
 C_d(\phi_L)\geq1-\frac1L+
       2\sum_{m=1}^{M}\operatorname{sinc}^2(LmK/d)
                                 |\tau_-(V_\varepsilon^{mK})|^2. \tag{CB.22}
\] All omitted Fourier coefficients are nonnegative. Take K to infinity at fixed M, use (CB.2), then let M grow. For \(t>0\) this gives \[
 \liminf C_d(\phi_L)\geq\Theta(t)-1/L,\qquad
 \Theta(t)=\sum_{m\in\mathbb Z}e^{-m^2t^2}>1.                  \tag{CB.23}
\] For t=0 the lower bound is infinite. If the pair measures converged on compact tests to a bounded density R(s) tending to 1 at infinity, then \(\int\phi_LR\to1\) as L grows. This contradicts (CB.23). A bounded subsequence of \(|\varepsilon|K^{1/4}/\sqrt N\) has a further subsequence with finite limiting \(t_{K}\), so (CB.4) follows. This is a necessary-condition argument, not an exchange of a growing test support with the semiclassical limit.

An equivalent moving-window consequence is \[
 \liminf\operatorname{Var}_{\theta}\!\bigl[
      \#\{l:\theta_l\in[\theta,\theta+L/d)\}\bigr]
       \geq L^2(\Theta(t)-1).
\] Indeed the finite variance is \(L^2(C_d(\phi_L)-1)+L\). The underlying positive Fourier test was already established in the earlier Floquet note; it is not claimed as a new result of this pass.

